% ************************** Thesis Organisation *****************************

\begin{organisation}
The first chapter is spent reviewing the fields of economics and neuroscience. We chronicle the development of theories of utility, before examining the evidence that the brain, and especially the dopaminergic midbrain circuitry, encodes utility. We also discuss the experimental task paradigms that animal research into utility has employed.

In Chapter 2, we compare these paradigms to a second-price auction developed by Becker et al., 1964 and explore the theory behind what makes these ‘Becker-DeGroot-Marshak’ (\nomenclature{BDM}{Becker-DeGroot-Marshak (auction)}) auctions appealing to behavioural economists and show theoretical results of a rational bidder. We conclude the chapter by looking at the growing neuroscientific literature of human economic decision-making using auction tasks and explain why we consider the BDM an appealing choice for investigating the function of neuronal circuits of reward. 

In Chapter 3 we detail a BDM auction task we developed and trained two rhesus macaques on. The behaviour of the primates on the task is explored, and the monkey's manipulation of tasks inputs and bidding behaviour is studied. The task behaviour is also compared to the same task limited to just one auctioned good and to more traditional binary choice tasks. We show convincing evidence that the animals are able to perform the BDM auction task \textit{as if} they are maximising some internal utility.

Chapter 4 concerns the electrophysiological recording of the dopaminergic cells of the midbrain of the two trained rhesus macaques. It begins with the single-cell recording procedure and the confirmation of the dopaminergic nature of the recorded cells. We then analyse the activity of dopaminergic cells as the animal responds to the auction task and makes their bid. We show that on a trial by trial level, dopaminergic midbrain cells in the animals encode various trial variables such as auction offer, the opposing computer bid, and whether or not the monkey won the trial. Crucially, we show that the activity of dopaminergic cells in the epoch where a good is offered is correlated with the subsequent monkey's bid.

Finally, we conclude by reviewing the results presented and the potential future directions that this work implies. We also briefly discuss the current state of primate/dopaminergic research and how neuroscientific studies can move beyond staid hypothesis.

\end{organisation}
